\chapter{Installation and running}

\eyebrow{installation}

\section{Option A: a pre-built release}

A native build for each platform is available on the project's \textbf{Releases}
page. Each release provides a self-contained archive per operating system, with
the model weights bundled and no build tools required:

\begin{center}
\begin{tabularx}{\linewidth}{@{}l X@{}}
\toprule
\textbf{Platform} & \textbf{Archive} \\
\midrule
macOS (universal) & \tele{carcass-macos.tar.gz}, contains the \tele{.app} \\
Windows (amd64)   & \tele{carcass-windows.zip}, contains the \tele{.exe} \\
Linux (amd64)     & \tele{carcass-linux.tar.gz}, contains the binary \\
\bottomrule
\end{tabularx}
\end{center}

The archive is unpacked and the application is run. The Python analysis features
still require a Python environment with the scientific libraries (see below); the
remaining features run without further setup. Releases are produced automatically
by continuous integration (Station~14).

\begin{note}[title={Building from source}]
To develop, modify, or build for a platform not covered by the releases, follow
Options~B and~C below.
\end{note}

\section{Prerequisites (build from source)}

\begin{center}
\begin{tabularx}{\linewidth}{@{}l l X@{}}
\toprule
\textbf{Component} & \textbf{Version} & \textbf{Needed for} \\
\midrule
Go            & 1.23+ & Building and running the core \\
Node.js       & 18+   & Building the front end \\
Wails CLI     & v2    & Development and native builds \\
Python        & 3.10+ & The inference / analysis sidecars \\
PyTorch, OpenCV, scikit-image, scipy, rembg & --- & Fat, conformation, live monitor \\
\bottomrule
\end{tabularx}
\end{center}

The core, capture and import features run with Go and Node alone. The model-based
features (fat analysis, conformation, live monitor) additionally require a Python
environment with the scientific libraries above.

\section{Option B: run from source}

\begin{lstlisting}
# install the Wails CLI
go install github.com/wailsapp/wails/v2/cmd/wails@latest

# python dependencies for the analysis sidecars
pip install torch opencv-python scikit-image scipy rembg onnxruntime
\end{lstlisting}

\subsection*{Model weights}

The trained weights and the ColorNaming lookup table are not shipped in the
source tree. Place them under \tele{model/weights/}:

\begin{center}
\begin{tabularx}{\linewidth}{@{}l X@{}}
\toprule
\textbf{File} & \textbf{Purpose} \\
\midrule
\tele{fat\_binary.pth}        & Fat segmentation (validated, IoU $\sim$0.92) \\
\tele{direct\_class.pth}      & Finishing grade (experimental) \\
\tele{eg\_regression.pth}     & Fat-thickness regression (experimental) \\
\tele{joost\_color\_naming.mat} & ColorNaming lookup table (required by all) \\
\bottomrule
\end{tabularx}
\end{center}

\section{Running in development}

\begin{lstlisting}
# hot-reload development (front end + Go rebuild on change)
wails dev
\end{lstlisting}

If the system \tele{python3} does not have the scientific libraries, the app can
be directed to one that does:

\begin{lstlisting}
export CARCASS_PYTHON=/path/to/python-with-torch
wails dev
\end{lstlisting}

\begin{note}[title=Camera and the WebView]
Webcam capture and the live monitor use the system camera through the WebView, so
they must run inside the app (via \tele{wails dev} or the compiled binary), not in
a plain browser. The bridge to the Go core only exists inside the app.
\end{note}

\section{Option C: build a native app}

\begin{lstlisting}
wails build -platform darwin/universal   # .app
wails build -platform windows/amd64      # .exe
wails build -platform linux/amd64        # binary
\end{lstlisting}

For the packaged app, the model weights are copied into the application's
resources so the analysis sidecars can find them; the helper script
\tele{scripts/build.sh} does this automatically after a build.
